\chapter*{Notation and conventions}
\addcontentsline{toc}{chapter}{Notation and conventions}
\markboth{Notation and conventions}{Notation and conventions}

\paragraph*{Valued fields.}
We write $(K,v_K)$ for a valued field, with valuations written additively, and use
\[
    \OK, \qquad \mK, \qquad k=\OK/\mK, \qquad
    \Gamma_K=v_K(K^\times)
\]
for its valuation ring, maximal ideal, residue field, and value group.
When a discrete valuation is normalized, our convention is $v_K(K^\times)=\ZZ$.
Completeness, discreteness, and assumptions on the residue field are stated where they are needed and are not imposed globally unless explicitly indicated.

\paragraph*{Curves and function fields.}
The letter $X$ usually denotes a smooth, projective, absolutely irreducible curve over $K$, called a \emph{nice curve} in Chapter~\ref{ch:models-of-curves-over-valued-fields}.
Its function field is denoted by $F_X=K(X)$.

\paragraph*{Models and fibres.}
A model of $X$ over $\OK$ is generally denoted by $\X$.
If $\eta$ and $s$ are respectively the generic and closed points of $\Spec\OK$, then
\[
    \X_\eta \quad\text{and}\quad \X_s
\]
denote the generic and special fibres.
Thus the book uses the French-style notation $\X_\eta$, not $\X_g$, for the generic fibre. % notation-check: allow documented deprecated form
For an irreducible component $Z$ of a special fibre, its generic point is denoted by $\xi_Z$, or simply $\xi$ when no confusion can arise; the associated geometric valuation is $v_Z$, and the finite set of component valuations is $V(\X)$.

\paragraph*{Base change.}
For a scheme $Y$ over a ring $R$ and a ring homomorphism $R\to S$, we write
\[
    Y_S \coloneqq Y\times_{\Spec R}\Spec S.
\]
After this definition, we use the subscript notation $Y_S$; in particular, base changes of $X$ and $\X$ are denoted by $X_L$ and $\X_{\OL}$.
For rings, modules, algebras, and sheaves, base change is denoted by $\otimes_R$.
When $L/K$ is finite, the normalized base change of $\X$ is denoted by
\[
    (\X_{\OL})^\sim.
\]
Function-field composita are written $F_X\mathbin{\cdot}L$.
